% Anatomy of an episode E=(z,M,e_{1:T},y), and of a candidate region inside it.
% Authored artwork (TikZ), included by both the ACM and the article build.
% Colours are declared in each wrapper preamble so the figure is class-neutral.
%
% Geometry is explicit (x=y=1mm, anchor=north west) for the same reason as
% architecture.tex: the event lane, the region bracket and the barrier strip must
% not collide at any scale. Natural width is ~162mm, which fits the article
% measure (165mm) and the two-column figure* span (178mm) without scaling.
\begin{tikzpicture}[
  x=1mm, y=1mm,
  font=\sffamily\scriptsize,
  gev/.style={
    draw=gacline!85, line width=0.45pt, rounded corners=1.2pt, fill=white,
    anchor=north west, align=center, inner sep=1.6pt,
    minimum width=16mm, minimum height=9mm, text width=13.6mm,
    font=\sffamily\tiny},
  mod/.style={gev, fill=gacblue!9,  draw=gacblue!85},
  tres/.style={gev, fill=gacteal!12, draw=gacteal!85},
  gside/.style={
    draw=gacline!70, line width=0.5pt, rounded corners=1.6pt, fill=gacgrey,
    anchor=north west, align=left, inner sep=2.6pt,
    minimum width=24mm, text width=21mm, font=\sffamily\tiny},
  barrier/.style={
    draw=gacred!65, line width=0.45pt, rounded corners=1.2pt, fill=gacred!8,
    anchor=north west, align=center, inner sep=1.6pt,
    minimum width=17mm, minimum height=6.4mm, text width=14.6mm,
    font=\sffamily\tiny},
  flow/.style={-{Latex[length=1.4mm,width=1.1mm]}, draw=gacline!85, line width=0.45pt},
  gtick/.style={draw=gacline!55, line width=0.4pt},
  gnote/.style={font=\sffamily\tiny, text=gacline, anchor=north west, align=left},
  ghd/.style={font=\sffamily\bfseries\tiny, text=gacblue, anchor=north west},
]

% ---------------------------------------------------------------- entry side
\node[gside] (z) at (0,0)
  {\textbf{entry state $z$}\\\code{issue\_number}: 4420};
\node[gside] (man) at (0,-12)
  {\textbf{manifest $M$}\\prompt $\cdot$ policy $\cdot$ guardrail\\tools $\cdot$ model $\cdot$ SDK\\
   tracer $\cdot$ entry contract\\effect catalog};

% ------------------------------------------------------------- ordered events
\node[mod]  (m1) at (30,   0) {model\\request};
\node[tres] (t1) at (48,   0) {tool result\\\code{record}};
\node[mod]  (m2) at (66,   0) {model\\request};
\node[tres] (t2) at (84,   0) {tool result\\\code{labels}};
\node[mod]  (m3) at (103,  0) {model\\request};
\node[tres] (t3) at (121,  0) {tool result\\\code{comments}};
\node[gnote]  (dots) at (99.5,-3.4) {$\cdots$};

\node[gside, minimum width=18mm, text width=15mm, fill=white, draw=gacline!70]
  (y) at (141,0) {\textbf{outcome $y$}\\classification\\+ excerpt};

\draw[flow] (z.east) -- (m1.west);
\draw[flow] (t3.east) -- (y.west);
\foreach \a/\b in {m1/t1, t1/m2, m2/t2, m3/t3} {\draw[flow] (\a.east) -- (\b.west);}

\node[ghd] at (30,4.8) {ordered events $e_{1:T}$ \normalfont\itshape (one recorded execution, not a prompt/answer pair)};

% ------------------------------------------------------- candidate region R
% Drawn as a plain bracket rather than a decoration: neither wrapper loads
% decorations.pathreplacing, and the figure must stay class-neutral.
\draw[gtick] (30,-11.4) -- (30,-13.4) -- (100,-13.4) -- (100,-11.4);
\node[gnote, text width=124mm] at (30,-14.2)
  {\textbf{candidate region $R=[a,b]$}: begins at a model-request boundary, ends after a tool
   result.  The deployable compiler additionally requires $a=0$, Eq.~\eqref{eq:prefix}.  The
   outcome $y$ is never inside $R$ --- \method removes intermediate control decisions, not the
   task result.};

% -------------------------------------------------------------- barrier strip
\node[barrier] (b1) at (22,-27) {approval};
\node[barrier] (b2) at (41,-27) {handoff};
\node[barrier] (b3) at (60,-27) {write};
\node[barrier] (b4) at (79,-27) {error};
\node[barrier] (b5) at (98,-27) {unknown\\effect};
\node[barrier] (b6) at (117,-27) {commit\\boundary};
\node[gnote, text width=25mm, text=gacred!85] at (135,-27.4)
  {any of these inside $R$ is a \textbf{barrier}, not a slow read};

\end{tikzpicture}
